\chapter*{Résumé}
\addcontentsline{toc}{chapter}{Résumé}

Le présent rapport décrit la conception et la mise en œuvre d'une solution décisionnelle (Business Intelligence) destinée à l'analyse et à l'optimisation du transport urbain de la ville de Dakar, réalisée pour le compte du CETUD (Conseil Exécutif des Transports Urbains de Dakar). Le projet s'appuie sur les données de l'Enquête Ménages-Déplacements (EMD) ainsi que sur des données de comptage de trafic, et couvre l'ensemble de la chaîne décisionnelle~: compréhension métier, modélisation multidimensionnelle des données en schéma en constellation (3 faits, 7 dimensions), mise en œuvre d'un processus d'extraction-transformation-chargement (ETL) avec Talend, implémentation d'un entrepôt de données sous PostgreSQL, restitution à travers des tableaux de bord interactifs développés avec Knowage BI, et enrichissement de la solution par des modèles de Machine Learning (segmentation des usagers, système de recommandation, détection d'anomalies de trafic et prédiction des zones à risque d'inaccessibilité). L'ensemble de la solution est exposé à travers une application web (React / FastAPI) destinée à la fois aux usagers du réseau et aux décideurs du CETUD. Ce travail répond à la problématique de la valorisation d'un grand volume de données hétérogènes afin de fournir une vision globale, fiable et exploitable des flux de transport, en vue d'améliorer la prise de décision et l'optimisation du réseau de transport urbain.

\textbf{Mots-clés~:} Business Intelligence, Data Warehouse, ETL, Talend, PostgreSQL, Knowage, schéma en constellation, Machine Learning, transport urbain, CETUD, Dakar.

\clearpage
